% Part: first-order-logic
% Chapter: beyond
% Section: many-sorted-logic

\documentclass[../../../include/open-logic-section]{subfiles}

\begin{document}

\olfileid[fr]{fol}{byd}{msl}

\olsection{Logique multi-sortée}

En logique du premier ordre, les variables et les quantificateurs
parcourent un unique !!{domain}. Il est pourtant souvent utile de
disposer de plusieurs !!{domain}s disjoints : on peut, par exemple,
vouloir un !!{domain} de nombres, un autre d'objets géométriques,
un autre de fonctions des nombres vers les nombres, un autre encore
de groupes abéliens, et ainsi de suite.

La logique multi-sortée fournit un tel cadre. On commence par une liste
de « sortes » : la sorte d'un objet indique le « !!{domain} »
auquel il est censé appartenir. On dispose alors de !!{variable}s et
de quantificateurs pour chaque sorte et, en général, d'une relation
d'!!{identity} pour chaque sorte. Les fonctions et les relations
sont elles aussi typées selon les sortes des objets qu'elles peuvent
recevoir comme arguments ; les fonctions ont aussi une sorte de résultat. Pour le reste, on conserve les règles
usuelles de la logique du premier ordre, en répétant les règles des
quantificateurs pour chaque sorte.

Par exemple, pour étudier les relations internationales, nous
pourrions choisir un langage comportant deux sortes d'objets : les
citoyens français et les citoyens allemands. Nous pourrions avoir une
relation unaire « boit du vin » pour les objets de la première sorte,
une autre, « mange de la saucisse », pour ceux de la seconde, et une
relation binaire « forme un couple marié multinational » dont le
premier argument est de la première sorte et le second de la seconde.
Si les variables $a$, $b$, $c$ parcourent les citoyens français et
$x$, $y$, $z$ les citoyens allemands, alors
\[
\lforall[a][\lforall[x][(\Atom{\Obj{MarriedTo}}{a,x} \lif
(\Atom{\Obj{DrinksWine}}{a} \lor \lnot \Atom{\Obj{EatsWurst}}{x}))]]
\]
affirme que, si une personne française est mariée à une personne
allemande, alors la première boit du vin ou la seconde ne mange pas
de saucisse.\footnote{Correction éditoriale : nous rétablissons dans l'affichage les arguments du quantificateur intérieur, mal délimités dans la source, ainsi que la parenthèse extérieure omise. Pour obtenir des sortes disjointes, on distingue ici les objets par leur étiquette de sorte ; une même personne ayant les deux nationalités peut ainsi avoir une copie dans chaque sorte.}

On peut traduire naturellement la logique multi-sortée en logique du
premier ordre : on réunit tous les objets de ses différents
!!{domain}s dans un seul !!{domain} du premier ordre, on utilise des
!!{predicate}s unaires pour distinguer les sortes et l'on relativise
les quantificateurs. Ainsi, le langage du premier ordre correspondant
à l'exemple précédent aurait, outre les autres relations décrites,
les !!{predicate}s unaires « $\Obj{German}$ » et
« $\Obj{French}$ », les contraintes de sorte sur les autres
symboles ayant été effacées. Un quantificateur portant sur une sorte
$\lforall[x][!A]$, où $x$ est !!a{variable} de la sorte des citoyens allemands,
se traduit par
\[
\lforall[x][(\Atom{\Obj{German}}{x} \lif !A)].
\]
Il faut ajouter des axiomes garantissant que les sortes sont
distinctes, par exemple
$\lforall[x][\lnot (\Atom{\Obj{German}}{x} \land
  \Atom{\Obj{French}}{x})]$, ainsi que des axiomes assurant que
« boit du vin » ne vaut que des objets qui satisfont le prédicat
$\Atom{\Obj{French}}{x}$, et ainsi de suite. Avec ces conventions
et ces axiomes, on montre sans difficulté que les !!{sentence}s de
la logique multi-sortée se traduisent en !!{sentence}s du premier
ordre, et ses !!{derivation}s en !!{derivation}s du premier ordre.
De même, les !!{structure}s multi-sortées se « traduisent » en
!!{structure}s du premier ordre correspondantes et réciproquement;
on obtient donc aussi un théorème de complétude pour la logique
multi-sortée.\footnote{Précision éditoriale : dans cet exemple à deux sortes, la correspondance avec l'union disjointe suppose aussi que chaque sorte soit non vide et que leur réunion soit le domaine. Les constantes doivent appartenir à leur sorte ; les fonctions doivent envoyer les tuples d'arguments des sortes prescrites dans leur sorte de résultat. Un existentiel se relativise par $\lexists[x][(\Atom{\Obj{German}}{x}\land !A)]$, avec une conjonction. Ces conditions complètent l'esquisse de la source.}

\end{document}
